%
% spec-opencode-integration.tex
%
% OpenCode-to-Ethos Session and Delegation Binding: A Z Specification
% A formal, harness-neutral model of the session/mission/delegation
% correlation problem tracked by bead ethos-5y4r (mirrored as points-hqi
% in punt-labs/points). Drafted BEFORE any OpenCode-side design or
% implementation work, per the mission brief: this document is the
% design step, not a retrospective account of code that already exists
% for OpenCode (none does yet).
%
% The Claude Code half of this model is read off real code, cited by
% file and line, exactly as spec-mission-lifecycle.tex and
% spec-writeset-admission.tex are: internal/process/tree.go,
% internal/session/{session,store,match}.go, internal/resolve/resolve.go,
% internal/mission/active.go, internal/hook/{pretooluse_dispatch,
% commit_trailers,audit_log}.go, internal/githook/{resolver,githook}.go.
%
% The OpenCode half is grounded in a companion research artifact,
% docs/prfaq-opencode-integration.md (a sibling worktree), which was
% sourced directly from anomalyco/opencode's dev branch (specs/v2/
% session.md, V2_HTTP_API_AUDIT.md, the CLI-plugin cli.mdx, and the
% migrate-v1.mdx hook-surface reference) via context7, across two
% evaluation rounds -- the second of which corrected a V1-vocabulary
% draft and located the context.data.session.{root,family,status}
% surface a round-1 search had correctly disclosed as unfound rather
% than guessed at. This agent did not independently fetch those
% primary sources (no context7/WebFetch tool was available in this
% session) and attributes every OpenCode-specific claim below to that
% companion document -- see \S\ref{sec:harness-instantiation}'s
% opening paragraph for the precise chain of evidence.
%
% Compile: pdflatex spec-opencode-integration.tex
%
% Must be .tex, not .md -- for the identical reason given in
% spec-mission-lifecycle.tex's own header comment and restated in
% spec-writeset-admission.tex's: this repo's /z-spec:test tooling globs
% docs/*.tex for the fuzz usepackage (Makefile's FUZZ_SPECS), and
% probcli's Z-to-B translation silently produces 0 states against any
% other extension. The mission contract naming this file's write_set as
% docs/spec-opencode-integration.md is corrected here, in the same PR,
% to the .tex path the toolchain actually requires -- see this
% document's own mission result for the explicit note.
%
\documentclass[a4paper,10pt,fleqn]{article}
\usepackage[margin=1in]{geometry}
\usepackage{fuzz}
\usepackage{booktabs}
\usepackage{tabularx}
\usepackage[colorlinks=true,linkcolor=blue,citecolor=blue,urlcolor=blue]{hyperref}

\newcommand{\fpath}[1]{\texttt{#1}}
\newcommand{\field}[1]{\texttt{#1}}

\begin{document}

\title{OpenCode-to-Ethos Session and Delegation Binding: A Z Specification}
\author{Mike S, for Punt Labs}
\date{October 2026}
\hypersetup{
  pdftitle={OpenCode-to-Ethos Session and Delegation Binding: A Z Specification},
  pdfauthor={Punt Labs},
  pdfsubject={Z Specification},
  pdfcreator={fuzz/probcli}
}
\maketitle
\tableofcontents
\newpage

% ===================================================================
\section{Scope}
\label{sec:scope}
% ===================================================================

Ethos today binds a Claude Code session to its identity, mission-claim,
and delegation state through a chain that is, in several places,
Claude-Code-specific by construction: a corroborated \field{CLAUDE\_PID}
(\fpath{internal/process/tree.go}), a PID-keyed pointer file
(\fpath{internal/session/store.go}'s \field{WriteCurrentSession}/
\field{ReadCurrentSession}), and a resolution chain
(\fpath{internal/resolve/resolve.go}'s \field{SessionID}) that walks a
process tree when no environment override is present. Bead
\field{ethos-5y4r} (mirrored as \field{points-hqi} in
\fpath{punt-labs/points}) asks for a harness-neutral generalization of
this chain that ALSO covers OpenCode 2.x, specified formally before any
design or implementation work proceeds.

This document models the state, invariants, transitions, and failure
modes of that generalized binding: which (harness, session) pairs are
active, which delegations are open and to which session and mission
each belongs, and how a commit's Mission/Delegation trailers are
resolved -- or, correctly, refused -- from that state. It does not model
the mission-contract status machine itself (that is
spec-mission-lifecycle.tex's job) or write-set path semantics (that is
spec-writeset-admission.tex's) -- both are cross-referenced where the
boundary matters, never re-derived.

\subsection{Non-goals}
\label{sec:nongoals}

\begin{itemize}
  \item \textbf{Circuit's PR-watch workflow state} -- the mutable PR
    head SHA tracked by the sibling Circuit/GitHub integration is
    explicitly a separate, GitHub-owned concern per the bead's own
    design note (``\emph{mutable PR head SHA remains GitHub/Circuit-owned}'').
    This document covers only the STABLE correlation metadata: harness
    session ID, Bead, Ethos mission, and delegation.
  \item \textbf{Changing Claude Code's existing behavior} -- every
    schema and operation below that touches the Claude Code side is
    read off code that already ships; \S\ref{sec:harness-instantiation}
    shows that code as ONE instantiation of this generalized model, not
    a replacement for it. A reviewer checking this document against
    \fpath{internal/resolve/resolve.go} should find the existing
    \field{SessionID} three-outcome contract unchanged.
  \item \textbf{The audit-log enrichment fallback chain's own known gap}
    -- \fpath{internal/hook/audit\_log.go}'s \field{buildAuditEntry}
    enriches a tool-call audit entry with \field{delegation\_id}/
    \field{contract\_id} from, in order, the tool input itself, then
    \field{DELEGATION\_ID}/\field{MISSION\_ID} env vars, with no
    sidecar fallback for a Tier A spawn (\fpath{audit\_log.go}:189--232).
    That gap is tracked separately as \field{ethos-f298} and is a
    property of audit-entry ENRICHMENT at tool-call time, not of the
    session/delegation BINDING this document models -- a different
    layer, not re-derived here.
  \item \textbf{True concurrent-writer interleaving.} Every operation
    below is a single atomic Z step, exactly the simplification
    spec-mission-lifecycle.tex's own \S~``Modelling choice: one mission
    at a time'' and spec-writeset-admission.tex's Validation section
    both disclose for their own state machines. The real mutual
    exclusion -- \field{AcquireDispatchPendingLock}
    (\fpath{internal/mission/active.go}:476--550), the per-session
    roster flock (\fpath{internal/session/store.go}'s
    \field{withLock}), the per-mission and per-delegation flocks -- is
    what makes the atomic-step assumption sound in the real system; it
    is not itself re-derived in Z. \S\ref{sec:failure-concurrent}
    states exactly what this document DOES check given that assumption,
    and what it does not.
  \item \textbf{OpenCode's own execution-recovery semantics.} The V2
    session contract commits a write-ahead claim when a process-local
    busy period starts, and an unclean OpenCode server death preserves
    that claim for the next startup to resume -- but, per the
    contract's own words, ``\emph{recovery is bounded per claimed
    execution but does not guarantee exactly-once provider requests or
    tool effects}'' (\fpath{specs/v2/session.md}, \S~``Execution Is
    Process-Local,'' as read and cited in this document's companion
    research artifact, \fpath{docs/prfaq-opencode-integration.md}).
    Whether a retried tool call runs exactly once is an OpenCode-side
    concern this document's correlation metadata is unaffected by --
    that metadata is a durable file, re-read fresh on every query, not
    an in-memory session object a crash could corrupt. The ACTUAL
    safeguard against a double-applied effect is the mission worker's
    own commit-per-step discipline (already required by this repo's
    mission-contract success criteria), not anything this integration
    adds or needs to model.
  \item \textbf{The exact wire shape of an OpenCode plugin payload.}
    \S\ref{sec:harness-instantiation} names the confirmed V2 hook
    surfaces and event types a plugin subscribes to
    (\field{Plugin.define}, \field{ctx.event.subscribe},
    \field{session.created}/\field{session.updated}) and the confirmed
    \field{context.data.session} lookup API
    (\field{root}/\field{family}/\field{status}), but does not model
    the full JSON shape of an event payload or every method on the
    OpenCode SDK client -- only the subset this document's operations
    actually consume.
  \item \textbf{Mission-contract status, round budget, result gates,
    write-set admission.} Already specified, respectively, by
    spec-mission-lifecycle.tex and spec-writeset-admission.tex. This
    document's \field{MISSIONID} is exactly that other document's
    mission identifier, treated here as an opaque given set.
\end{itemize}

\subsection{Precedent}
\label{sec:precedent}

Binding one tool's identity or audit state across more than one
runtime is not a new problem class, and this design does not pretend
to invent its shape from nothing. Four precedents, surveyed in detail
in this document's companion research artifact
(\fpath{docs/prfaq-opencode-integration.md}, \S~``Exemplars''), bear
directly on the invariants below: W3C Trace Context's
\field{TRACEPARENT}/\field{TRACESTATE} header pair, which Claude Code
itself already reads, shows that a single correlation identifier
travelling as inert metadata works across runtimes sharing nothing
else -- exactly the shape of \field{MISSIONID}/\field{DELEGATIONID}
this document moves between harnesses; SLSA build provenance's
immutable-provenance-vs-mutable-build-state split is the formal
ancestor of this document's own stable-correlation-metadata-vs-Circuit's
-mutable-PR-SHA non-goal boundary (\S\ref{sec:nongoals}); the in-house
\field{model-provider}/Bifrost tool already solves the sibling
problem -- one control surface, several harnesses -- by writing into
each harness's own native extension point rather than inventing a
cross-harness daemon, the same design this document's
\S\ref{sec:harness-instantiation} follows for OpenCode; and, closest
to home and most directly load-bearing for
\field{FailClosedOnAmbiguity} (\S\ref{sec:theorems}), ethos's own
existing \field{Mission:}/\field{Delegation:} trailer mechanism has
ALREADY failed silently once in production -- a session-binding bug
meant the commit-msg hook ``stamped nothing'' the first time a
multi-session scenario exercised it, per this repo's own
\fpath{docs/MODERNIZATION\_CASE\_STUDY.tex} as quoted in the companion
research artifact's \S~``Exemplars.'' This document did not
independently re-read that case-study file (it does not exist in this
author's own worktree checkout at the time of writing); the claim is
attributed to the companion artifact, not independently verified
here. \field{FailClosedOnAmbiguity} exists precisely so this
integration does not discover the same failure mode a second time, in
a second harness, before fixing it formally the first time.

\subsection{Modelling choice: one global registry, not per-harness
state machines}
\label{sec:modelling-choice}

A tempting alternative design gives Claude Code and OpenCode each their
own state schema and asks whether the two satisfy ``the same''
invariants by analogy. That throws away exactly the property worth
proving: that a SINGLE state machine, with no branch on which harness a
session belongs to, accounts for both. \S\ref{sec:state} defines one
schema, \field{Registry}, whose session and delegation components carry
a \field{HARNESS} tag purely as data -- no operation's guard or effect
ever inspects it. \field{HarnessNeutrality}
(\S\ref{sec:theorems}) states this precisely: it is a claim checkable by
inspection of every operation's predicate, not an empirical property
established by running the model twice under different parameters.

% ===================================================================
\section{Basic Types}
\label{sec:basic}
% ===================================================================

Three given sets carry no internal structure this specification needs:
a session identifier (a corroborated \field{CLAUDE\_PID}, an
\field{ETHOS\_SESSION} value, or -- per
\S\ref{sec:harness-instantiation} -- an OpenCode session ID), a mission
identifier (spec-mission-lifecycle.tex's own \field{MISSIONID}, opaque
here), and a delegation identifier (spec-mission-lifecycle.tex's
\field{blockingDelegationCount} counts these; this document tracks each
one's own session and mission binding instead of only a count, because
the binding -- not merely the count -- is exactly what this bead is
about).

\begin{zed}
[SESSIONID, MISSIONID, DELEGATIONID]
\end{zed}

% ===================================================================
\section{Free Types}
\label{sec:freetypes}
% ===================================================================

\subsection{ZBOOL}

Carried over from both existing specifications, for the identical
reason: \fuzz{} has no built-in boolean, and \field{BOOL}/\field{true}/
\field{false} collide with B keywords under ProB's translation.

\begin{zed}
ZBOOL ::= ztrue | zfalse
\end{zed}

\subsection{HARNESS}

The two harnesses this bead names. A third would be a third
constructor, not a design change -- nothing below branches on which
value a session carries (\S\ref{sec:modelling-choice}).

\begin{zed}
HARNESS ::= hClaude | hOpenCode
\end{zed}

\subsection{TIER}

The two delegation tiers \fpath{internal/mission/delegation.go}:19--20
names directly: \field{TierA} (an ad hoc spawn, no mission) and
\field{TierB} (spawned under an explicit mission).

\begin{zed}
TIER ::= tierA | tierB
\end{zed}

\subsection{ORIGIN}

How a Tier B delegation's mission came to be known, generalizing
\fpath{internal/mission/active.go}'s \field{BindOrigin*}/\field{BoundVia*}
constants (lines 100--113, 59--62) to the dimension this document
checks: was the binding stated outright, dispatched by a leader for
someone else, inherited down a \field{parent\_delegation} CHAIN (a
per-delegation pointer, the Claude Code Tier B inheritance mechanism),
or propagated by a root/family LOOKUP on the forking session's own
hierarchy (the OpenCode mechanism, \S\ref{sec:opforkpropagate}) --
these are two structurally different evidencing mechanisms for
``this delegation's mission binding is traceable, not asserted,'' kept
as two distinct constructors rather than collapsed into one, because
\S\ref{sec:delegation-schema}'s guard on each is different: one
requires a stored \field{DELEGLINK} pointer, the other requires
none, because the evidence lives in \field{SESSION.root} instead (see
\field{TraceabilityViaRootSuffices}, \S\ref{sec:theorems}).
\field{oNone} is reserved for Tier A, which has no mission and hence no
origin to record.

\begin{zed}
ORIGIN ::= oNone | oClaim | oDispatch | oInherited | oForkPropagated
\end{zed}

\subsection{PARENTLINK}

A session's fork provenance -- present (\field{forkedFrom(p)}) when
this session was created by forking an existing session \field{p},
absent (\field{noParent}) for a root session. Retained purely as
descriptive data on a \field{SESSION} record
(\S\ref{sec:session-schema}): unlike the ORIGINAL draft of this
document, no operation below ever walks this link to RESOLVE anything
-- the OpenCode research confirms lineage is looked up directly via
\field{context.data.session.root}/\field{.family}
(\S\ref{sec:harness-instantiation}), which this document models as the
\field{root}/\field{family} FIELDS on \field{SESSION} itself
(\S\ref{sec:session-schema}), not as a chain an operation has to
traverse at read time. \field{PARENTLINK} survives only as the
immediate-parent provenance a \field{ForkSession} event records, kept
separate from \field{root}/\field{family} because ``who directly forked
me'' and ``who is the lineage root / who else shares it'' are
genuinely different questions \S\ref{sec:opforksession} needs to ask
and answer independently.

\begin{zed}
PARENTLINK ::= noParent | forkedFrom \ldata SESSIONID \rdata
\end{zed}

\subsection{MISSIONLINK, DELEGLINK}

The same optional-value idiom, for a delegation's mission (absent on
Tier A) and a delegation's inherited parent delegation (absent except
under \field{oInherited}).

\begin{zed}
MISSIONLINK ::= noMission | missionOf \ldata MISSIONID \rdata
\end{zed}

\begin{zed}
DELEGLINK ::= noDeleg | delegOf \ldata DELEGATIONID \rdata
\end{zed}

% ===================================================================
\section{Global Constants}
\label{sec:constants}
% ===================================================================

Bounding the number of concurrently tracked sessions and delegations is
what lets probcli enumerate the state space at all -- the identical
economy argument spec-mission-lifecycle.tex's \field{maxRounds}/
\field{maxDelegations} and spec-writeset-admission.tex's \field{maxOpen}
make for their own bounds. Neither constant has a real-code analogue:
the actual \field{Registry} (a session roster directory plus a mission
tree) holds arbitrarily many of each.

\begin{axdef}
maxSessions : \nat
\where
maxSessions = 2
\end{axdef}

\begin{axdef}
maxDelegations : \nat
\where
maxDelegations = 2
\end{axdef}

\field{RestartService} (\S\ref{sec:oprestart}) increments
\field{Registry.epoch} unconditionally on every firing, so without an
explicit bound the epoch component is unbounded in exactly the way
z-spec's own ProB-compatibility guidance warns against -- ProB cannot
enumerate a state space with an ever-growing natural number in it, and
a model-check run against the unbounded form times out rather than
reporting anything. \field{maxEpoch} bounds it purely for this reason;
the real system has no such ceiling (a long-lived \field{ethos serve}
process restarts an unbounded number of times over its lifetime).

\begin{axdef}
maxEpoch : \nat
\where
maxEpoch = 2
\end{axdef}

One small partial function recovers the wrapped value from a
\field{DELEGLINK}, exactly as spec-writeset-admission.tex's
\field{unwrap} does for \field{RAWSEG}: undefined on the ``absent''
constructor, total and faithful on the ``present'' one.

\begin{axdef}
delegID : DELEGLINK \pfun DELEGATIONID
\where
\dom delegID = \ran delegOf \\
\forall d : DELEGATIONID @ delegID(delegOf(d)) = d
\end{axdef}

A second such function recovers the wrapped value from a
\field{PARENTLINK}: unlike the first draft of this document, it is used
only to STATE the lineage invariant below (\S\ref{sec:registry}), never
to RESOLVE a chain of pointers at query time -- \S\ref{sec:registry}'s
\field{lineageRoot} is the maintained projection that makes repeated
chain-walking unnecessary, matching how OpenCode's own
\field{context.data.session.root}/\field{.family} are confirmed,
already-maintained lookups, not something a plugin has to derive itself
on every call (\S\ref{sec:harness-instantiation}).

\begin{axdef}
parentID : PARENTLINK \pfun SESSIONID
\where
\dom parentID = \ran forkedFrom \\
\forall s : SESSIONID @ parentID(forkedFrom(s)) = s
\end{axdef}

% ===================================================================
\section{State}
\label{sec:state}
% ===================================================================

\subsection{SESSION}
\label{sec:session-schema}

The generalized harness session identity -- the minimal common shape a
Claude Code session and an OpenCode session both instantiate. Five
fields: which harness, this session's immediate fork provenance (if
any), the lineage ROOT session (every session is a member of exactly
one lineage, rooted at itself if it was never forked from anything),
the liveness epoch at which this record was last corroborated as real,
and whether it is currently recorded as alive at all.

\begin{schema}{SESSION}
harness : HARNESS \\
parent : PARENTLINK \\
root : SESSIONID \\
bornEpoch : \nat \\
alive : ZBOOL
\where
bornEpoch \leq maxEpoch
\end{schema}

\field{root} is this document's Z rendering of OpenCode's own
\field{context.data.session.root(sessionID)} lookup
(\S\ref{sec:harness-instantiation}) -- a direct, maintained hierarchy
projection the CLI-plugin data-cache API exposes, independent of
whatever fields a \field{session.created} event payload happens to
carry. A session's FAMILY -- OpenCode's
\field{context.data.session.family(sessionID)} -- is deliberately NOT
a second stored field on \field{SESSION}: it is the DERIVED set of
every currently-rostered session sharing the same \field{root}
(\field{Registry}'s own \field{family} function,
\S\ref{sec:registry}). Storing family as data on every member would
require rewriting every other family member's record each time a new
fork joins -- an operation with no natural single-key Z rendering and
no counterpart in the real API either, which exposes \field{family} as
a LOOKUP, not a field a caller sets. Deriving it from \field{root}
equality is the one representational change this document makes
between its first and second draft, driven directly by confirming
OpenCode's hierarchy API shape; \field{ForkSession}
(\S\ref{sec:opforksession}) therefore needs to write only the ONE new
session's own record, never a bulk update over its lineage.

\field{bornEpoch} together with \field{Registry}'s own \field{epoch}
component (\S\ref{sec:registry}) is the generalization of
\field{ClaudePIDFromEnvCorroborated}'s live-ancestor check
(\fpath{internal/process/tree.go}:52--86): a session is usable only
when its OWN liveness anchor has been corroborated against CURRENT
reality, not merely recorded once and cached. \field{FindClaudePID}'s
own doc comment states the property this generalizes directly:
``\emph{resolved fresh on every call -- never cached. A long-lived
process (\field{ethos serve}) outlives any single session it has
observed; a process-lifetime cache... would keep answering with the
first session's PID forever}'' (\fpath{tree.go}:18--22, DES-074). A
session record is never mutated to look alive by assertion; only
\field{CorroborateSession} (\S\ref{sec:opcorroborate}) can re-stamp
\field{bornEpoch}, and only by re-asserting every other field unchanged.
\S\ref{sec:harness-instantiation} shows this abstract pair
INSTANTIATED two different ways: a PID-ancestry re-walk for Claude
Code, and the confirmed, cheaper
\field{context.data.session.status(sessionID)} projection for
OpenCode -- the status primitive the companion research names as
materially better than reconstructing liveness from the write-ahead
execution-claim recovery language (\S\ref{sec:nongoals}'s own
disclosed exclusion of that recovery machinery).

\subsection{DELEGATION}
\label{sec:delegation-schema}

One delegation's binding: which session opened it, which mission (if
any) it was opened under, which tier, how its mission binding was
established, and -- for an inherited binding -- which ancestor
delegation licensed the inheritance.

\begin{schema}{DELEGATION}
session : SESSIONID \\
mission : MISSIONLINK \\
tier : TIER \\
origin : ORIGIN \\
parentDelegation : DELEGLINK \\
bornEpoch : \nat
\where
bornEpoch \leq maxEpoch \\
tier = tierA \implies
  (mission = noMission \land origin = oNone \land parentDelegation = noDeleg) \\
tier = tierB \implies (mission \neq noMission \land origin \neq oNone) \\
origin = oInherited \iff parentDelegation \neq noDeleg
\end{schema}

$\field{bornEpoch} \leq \field{maxEpoch}$ is the same ProB-animatability bound
\S\ref{sec:registry} adds to \field{Registry.epoch}, needed here for
the identical reason. The other three conjuncts are, respectively: a
Tier A delegation
(\fpath{internal/mission/delegation.go}:19, \fpath{internal/hook/
pretooluse\_dispatch.go}'s \field{dispatchTierA}) carries no mission and
needs no origin to explain, because there is nothing to attribute; a
Tier B delegation always names a mission, by definition of the tier;
and -- the formal statement of the bead's own ``a fork of a session
never silently inherits its parent's mission binding without an
explicit, traceable link'' criterion -- \field{origin = oInherited}
IF AND ONLY IF a non-empty \field{parentDelegation} pointer is carried.
Strengthened from the first draft's one-directional implication to a
biconditional once \field{oForkPropagated} was added
(\S\ref{sec:opforkpropagate}): a fork-propagated delegation's own
traceability witness is a shared \field{family}, never a
\field{parentDelegation} pointer, so the type now makes it
UNREPRESENTABLE for any non-\field{oInherited} origin -- including the
new one -- to carry that pointer at all, not merely unconventional.
There is no value of \field{ORIGIN} that lets a delegation claim
``inherited'' while leaving the chain unrecorded, nor one that lets a
delegation carry an inheritance pointer while claiming a DIFFERENT
origin: the type itself makes a silent or a misattributed inheritance
unrepresentable, not merely disallowed by a guard that a caller could
omit.

\subsection{Registry}
\label{sec:registry}

The one global state: every tracked session, every tracked delegation,
every mission claim, every session's delegation-binding sidecar, and
the current liveness epoch.

\begin{schema}{Registry}
sessions : SESSIONID \pfun SESSION \\
delegations : DELEGATIONID \pfun DELEGATION \\
claims : SESSIONID \pfun MISSIONID \\
delegationOf : SESSIONID \pfun DELEGATIONID \\
epoch : \nat
\where
epoch \leq maxEpoch \\
\# (\dom sessions) \leq maxSessions \\
\# (\dom delegations) \leq maxDelegations \\
\dom claims \subseteq \dom sessions \\
\dom delegationOf \subseteq \dom sessions \\
\forall s : \dom delegationOf @
  delegationOf(s) \in \dom delegations \land
  (delegations(delegationOf(s))).session = s \\
\forall s : \dom sessions |
    (sessions(s)).parent = noParent @
  (sessions(s)).root = s \\
\forall s : \dom sessions |
    (sessions(s)).parent \neq noParent \land
    parentID((sessions(s)).parent) \in \dom sessions @
  (sessions(s)).root = (sessions(parentID((sessions(s)).parent))).root
\end{schema}

\field{family(s)}, the derived per-session lookup OpenCode's own
\field{context.data.session.family(sessionID)} names
(\S\ref{sec:session-schema}), is NOT declared as a free-standing
global function here, deliberately: spec-writeset-admission.tex's own
\S~``Modelling choice: schemas, not free functions, over sequences''
already identifies the reason -- declaring any global constant whose
type mentions a function or relation over an unbounded carrier makes
\texttt{SETUP\_CONSTANTS} fail for the WHOLE machine, even when every
operation that actually fires never references it, because $\field{SESSIONID}
\pfun \field{SESSION}$ is unbounded as a Z type regardless of how
tightly \field{maxSessions} constrains it in practice. \field{family(s)}
is therefore written inline, at its one call site
(\S\ref{sec:opforkpropagate}), as the set comprehension
$\{ t : \dom sessions \mid (sessions(t)).root = (sessions(s)).root \}$
directly against whichever \field{Registry} binding is in scope -- the
identical VALUE the generic notation \field{family(sessions)} names
informally throughout this document's prose, with no schema-level
declaration for ProB to choke on.

Eight predicates, five load-bearing:

\begin{itemize}
  \item $\field{epoch} \leq \field{maxEpoch}$ and the two cardinality bounds are
    ProB animatability conditions, with no code counterpart -- the
    identical role spec-mission-lifecycle.tex's
    \field{blockingDelegationCount} $\leq$ \field{maxDelegations} and
    spec-writeset-admission.tex's \field{Store} invariant play for
    their own bounds. \field{epoch}'s bound is forced by
    \field{RestartService} (\S\ref{sec:oprestart}), the one operation
    whose effect is an unconditional increment; without it the state
    space is unbounded and probcli times out rather than reporting
    anything, exactly the failure this document's own first model-check
    attempt hit before \field{maxEpoch} was added.
  \item \field{claims} and \field{delegationOf} are both scoped to
    LIVE-ROSTER sessions ($\dom \field{claims}$ $\subseteq$
    $\dom \field{sessions}$) -- these are the two per-session SIDECARS
    (\field{active-mission}, \field{delegation-binding};
    \fpath{internal/mission/active.go}:31--49, 342--347), ephemeral
    state that lives beside a session's roster entry and is cleared
    when the roster entry goes (\field{EndSession},
    \S\ref{sec:opendsession}).
  \item \field{delegationOf}'s own defining conjunct -- a per-session
    sidecar always names a REAL delegation actually bound to THAT
    session -- is the formal shape of
    \field{WriteCommitTrailers}'s own same-mission cross-check
    (\fpath{internal/hook/commit\_trailers.go}:65: ``\emph{The
    delegation is emitted only when its binding names the same
    mission}''), restated as a standing state invariant rather than a
    read-time guard: the sidecar and the record it names can never
    drift apart, by construction, so \field{ResolveCommitAttribution}
    (\S\ref{sec:opresolve}) never needs to reconcile a mismatch -- there
    is none reachable.
  \item The last two conjuncts together are this redesign's
    \field{root} CONSISTENCY property, in place of the earlier draft's
    single dangling-\field{PARENTLINK} sanity bound: a session with no
    recorded fork-parent is its OWN lineage root, and a session WITH a
    recorded fork-parent, WHILE that parent is still currently
    rostered, agrees with it on which root they share. The guard is
    deliberately conditioned on the parent STILL being in
    $\dom \field{sessions}$ -- a fork's own \field{root} is stamped once,
    at \field{ForkSession} time, and never re-derived afterward
    (\S\ref{sec:opforksession}), so a parent ending its session later
    (\field{EndSession}, \S\ref{sec:opendsession}) must not retroactively
    falsify this invariant for every child it ever forked, exactly the
    same durability argument this document's own note on
    \field{(delegations(d)).session} already makes below. Deriving
    \field{family} from \field{root} EQUALITY (rather than storing it
    per-session) is what makes the earlier draft's separate
    ``no two family-mates disagree on root'' conjunct unnecessary here:
    two sessions sharing a root are family by \field{family}'s own
    definition, not by a second fact that could drift out of sync with
    the first.
\end{itemize}

Deliberately ABSENT from this schema: a conjunct requiring every
delegation's \field{session} field to name a currently-rostered
session. A delegation record is DURABLE audit data -- it lives under
the mission's own tree, independent of the session roster -- while a
session's roster entry is EPHEMERAL, removed on \field{EndSession} or
purge. Requiring $\field{(delegations(d)).session} \in \dom \field{sessions}$ for
every \field{d} would make \field{EndSession} (\S\ref{sec:opendsession})
retroactively falsify the state invariant for every delegation that
session ever opened -- exactly backwards, since the real
\fpath{internal/session/store.go}'s \field{Store.Delete} deletes a
roster entry while leaving the delegation tree under
\fpath{.punt-labs/ethos/missions/} entirely untouched. This omission is
a checked design decision, not an oversight: it is what lets
\field{EndSession} below drop a session from $\dom \field{sessions}$
without disturbing a single delegation record that names it.

% ===================================================================
\section{Initialization}
\label{sec:init}
% ===================================================================

\begin{schema}{Init}
Registry'
\where
sessions' = \emptyset \\
delegations' = \emptyset \\
claims' = \emptyset \\
delegationOf' = \emptyset \\
epoch' = 0
\end{schema}

% ===================================================================
\section{Operations}
\label{sec:ops}
% ===================================================================

Every schema below that constructs a new \field{SESSION} or
\field{DELEGATION} value takes it as an INPUT parameter
(\field{rec?}), constrained by equations, rather than building one from
field-literals inline -- the same discipline
spec-writeset-admission.tex's \field{AdmitMission} applies to
\field{entry? : ENTRY}. It avoids any dependence on a schema-binding
display construct and keeps every operation's predicate a flat
conjunction, matching house style.

\subsection{JoinRootSession}
\label{sec:opjoin}

Models a new, UNFORKED harness session entering the registry -- the
real \fpath{internal/session/store.go}'s \field{Store.CreateInCheckout}
for Claude Code (every Claude Code session is a root session per
\S\ref{sec:harness-instantiation}), and, for OpenCode, the
\field{session.created} event firing for a session with no fork
ancestry. A root session is its own lineage: \field{root} names
itself, so \field{family} (\S\ref{sec:registry}'s derived lookup)
evaluates to the singleton containing only itself the instant it joins
-- \field{ForkSession} (\S\ref{sec:opforksession}) is the only
operation whose effect ever makes some OTHER session's \field{family}
grow beyond one member.

\begin{schema}{JoinRootSession}
\Delta Registry \\
sid? : SESSIONID \\
rec? : SESSION
\where
sid? \notin \dom sessions \\
\# (\dom sessions) < maxSessions \\
rec?.parent = noParent \\
rec?.root = sid? \\
rec?.bornEpoch = epoch \\
rec?.alive = ztrue \\
sessions' = sessions \oplus \{ sid? \mapsto rec? \} \\
delegations' = delegations \\
claims' = claims \\
delegationOf' = delegationOf \\
epoch' = epoch
\end{schema}

\subsection{ForkSession}
\label{sec:opforksession}

Models a session FORKING an existing one -- OpenCode's
\field{POST /api/session/\{sessionID\}/fork} (operation ID
\field{session.fork}), confirmed by the companion research as a
first-class, routinely-exercised V2 primitive, not a corner case
(\S\ref{sec:harness-instantiation}). The forked child receives its OWN
session ID distinct from its parent's -- nothing about the fork
endpoint's contract ties the two identities together beyond transcript
and instruction-value copying, per the companion research's direct
reading of \fpath{specs/v2/session.md}'s own ``Instructions Are Value
Deltas'' section. The child's \field{root} is inherited from the PARENT
(a fork of a fork still belongs to the ORIGINAL lineage, never starts a
new one); because \field{family} is DERIVED from \field{root} equality
(\S\ref{sec:registry}), stamping \field{root} correctly on the one new
record is all that is needed for the new child to appear in its whole
lineage's \field{family} and for every existing member's own
\field{family} to pick the newcomer up automatically, the next time
either is queried -- modelling OpenCode's own
\field{context.data.session.family(sessionID)} as a maintained,
symmetric projection with no corresponding WRITE this operation has to
perform on any session other than the one it creates.

\begin{schema}{ForkSession}
\Delta Registry \\
sid? : SESSIONID \\
parentSid? : SESSIONID \\
rec? : SESSION
\where
sid? \notin \dom sessions \\
\# (\dom sessions) < maxSessions \\
parentSid? \in \dom sessions \\
rec?.parent = forkedFrom(parentSid?) \\
rec?.root = (sessions(parentSid?)).root \\
rec?.bornEpoch = epoch \\
rec?.alive = ztrue \\
sessions' = sessions \oplus \{ sid? \mapsto rec? \} \\
delegations' = delegations \\
claims' = claims \\
delegationOf' = delegationOf \\
epoch' = epoch
\end{schema}

The load-bearing conjunct for the bead's fork criterion is the one that
is ABSENT: \field{claims' = claims}. A freshly forked session enters
the registry with NO mission claim of its own at all -- the ONLY ways a
session's \field{claims} entry is ever populated are
\field{ClaimMission} (\S\ref{sec:opclaim}, an explicit operator act)
and \field{AllocateDelegationTierB}'s \field{oInherited} or
\field{oForkPropagated} cases (\S\ref{sec:opallocateb},
\S\ref{sec:opforkpropagate}, both of which require an explicit,
checkable traceability witness -- a \field{parentDelegation} link or a
shared \field{family}, respectively -- never a bare assertion).
\field{ForkNeverInheritsClaim} (\S\ref{sec:theorems}) names this
precisely, and is SHARPENED rather than weakened by \field{ForkSession}
being a frequent, first-class operation: the more often a fork fires,
the more this guarantee is actually exercised, not merely
theoretically available.

\subsection{ClaimMission}
\label{sec:opclaim}

Models \field{ethos mission claim} -- the operator declaring ``I am
working on this mission'' -- via \field{WriteActiveMission}
(\fpath{internal/mission/active.go}:223--229).

\begin{schema}{ClaimMission}
\Delta Registry \\
sid? : SESSIONID \\
mission? : MISSIONID
\where
sid? \in \dom sessions \\
(sessions(sid?)).alive = ztrue \\
(sessions(sid?)).bornEpoch = epoch \\
claims' = claims \oplus \{ sid? \mapsto mission? \} \\
sessions' = sessions \\
delegations' = delegations \\
delegationOf' = delegationOf \\
epoch' = epoch
\end{schema}

\subsection{ReleaseClaim}

Models \field{ethos mission release} clearing the active-mission
sidecar (\field{ClearActiveMission}, \fpath{active.go}:414--423).

\begin{schema}{ReleaseClaim}
\Delta Registry \\
sid? : SESSIONID
\where
sid? \in \dom claims \\
claims' = \{ sid? \} \ndres claims \\
sessions' = sessions \\
delegations' = delegations \\
delegationOf' = delegationOf \\
epoch' = epoch
\end{schema}

\subsection{AllocateDelegationTierA}

Models \field{dispatchTierA} (\fpath{internal/hook/
pretooluse\_dispatch.go}:657--700): an ad hoc spawn with no mission,
allocated a fresh \field{delegation\_id} and attributed to the issuing
session alone.

\begin{schema}{AllocateDelegationTierA}
\Delta Registry \\
sid? : SESSIONID \\
id? : DELEGATIONID \\
rec? : DELEGATION
\where
sid? \in \dom sessions \\
(sessions(sid?)).alive = ztrue \\
(sessions(sid?)).bornEpoch = epoch \\
id? \notin \dom delegations \\
\# (\dom delegations) < maxDelegations \\
rec?.session = sid? \\
rec?.mission = noMission \\
rec?.tier = tierA \\
rec?.origin = oNone \\
rec?.parentDelegation = noDeleg \\
rec?.bornEpoch = epoch \\
delegations' = delegations \oplus \{ id? \mapsto rec? \} \\
sessions' = sessions \\
claims' = claims \\
delegationOf' = delegationOf \\
epoch' = epoch
\end{schema}

\subsection{AllocateDelegationTierB}
\label{sec:opallocateb}

Models \field{dispatchTierB} (\fpath{pretooluse\_dispatch.go}:711--912):
a spawn attributed to an explicit mission, by one of the three routes
\fpath{active.go}'s \field{BoundVia*} constants name --
\field{BoundViaMissionIDEnv}/\field{BoundViaSidecarClaim} both collapse
to \field{oClaim} or \field{oDispatch} at this document's granularity
(the distinction between an explicit env var and a sidecar claim is a
mechanism-of-discovery detail with no bearing on any invariant checked
here -- the same economy spec-writeset-admission.tex applies when it
elides \field{segmentMatches}' per-segment classification), while
\field{BoundViaInherited} maps to \field{oInherited} exactly.

\begin{schema}{AllocateDelegationTierB}
\Delta Registry \\
sid? : SESSIONID \\
id? : DELEGATIONID \\
rec? : DELEGATION
\where
sid? \in \dom sessions \\
(sessions(sid?)).alive = ztrue \\
(sessions(sid?)).bornEpoch = epoch \\
id? \notin \dom delegations \\
\# (\dom delegations) < maxDelegations \\
rec?.session = sid? \\
rec?.tier = tierB \\
rec?.mission \neq noMission \\
rec?.origin \in \{ oClaim, oDispatch, oInherited \} \\
rec?.origin = oInherited \implies
  ( rec?.parentDelegation \neq noDeleg \land
    delegID(rec?.parentDelegation) \in \dom delegations \land
    (delegations(delegID(rec?.parentDelegation))).mission = rec?.mission ) \\
rec?.origin \neq oInherited \implies rec?.parentDelegation = noDeleg \\
rec?.bornEpoch = epoch \\
delegations' = delegations \oplus \{ id? \mapsto rec? \} \\
sessions' = sessions \\
claims' = claims \\
delegationOf' = delegationOf \\
epoch' = epoch
\end{schema}

The \field{oInherited} guard is the formal mirror of
\field{enforceDelegationDepth}'s own parent-chain walk
(\fpath{pretooluse\_dispatch.go}:912--964): the ancestor named by
\field{parentDelegation} must already be on record, and must name the
SAME mission the child inherits -- inheritance never silently invents a
mission the ancestor did not itself carry. The guard is narrowed from
the first draft's $\field{rec?.origin} \neq \field{oNone}$ to an explicit
three-way enumeration EXCLUDING \field{oForkPropagated}: a fork's own
delegation round is allocated by the DEDICATED
\field{AllocateDelegationForkPropagated} operation below
(\S\ref{sec:opforkpropagate}), never by this one, because the two have
different traceability witnesses to check (a \field{parentDelegation}
pointer versus a shared \field{family}) and conflating them into one
schema's guard would let either witness silently satisfy the other.

\subsection{AllocateDelegationForkPropagated}
\label{sec:opforkpropagate}

Models the companion research's own recommended plugin behavior for a
forked OpenCode session: on the child's \field{session.created}/
equivalent event firing with a new Session ID and no existing claim,
the plugin resolves lineage via \field{context.data.session.root}/
\field{.family} and, finding an active mission claim on some member of
that family, opens a NEW DELEGATION ROUND for the forked session under
the SAME mission -- ``ethos's mission model already supports multiple
delegations per mission, so this is an existing primitive applied to a
new trigger, not new machinery''
(\fpath{docs/prfaq-opencode-integration.md}, \S~``What happens if
OpenCode forks a session mid-mission?''). Given its OWN dedicated
operation, distinct from \field{AllocateDelegationTierB}, because V2
fork is frequent and first-class, not a corner case to fold into the
generic inheritance path.

\begin{schema}{AllocateDelegationForkPropagated}
\Delta Registry \\
sid? : SESSIONID \\
id? : DELEGATIONID \\
witness? : SESSIONID \\
rec? : DELEGATION
\where
sid? \in \dom sessions \\
(sessions(sid?)).alive = ztrue \\
(sessions(sid?)).bornEpoch = epoch \\
sid? \notin \dom claims \\
witness? \in \dom sessions \\
(sessions(witness?)).root = (sessions(sid?)).root \\
witness? \in \dom claims \\
id? \notin \dom delegations \\
\# (\dom delegations) < maxDelegations \\
rec?.session = sid? \\
rec?.tier = tierB \\
rec?.mission = missionOf(claims(witness?)) \\
rec?.origin = oForkPropagated \\
rec?.parentDelegation = noDeleg \\
rec?.bornEpoch = epoch \\
delegations' = delegations \oplus \{ id? \mapsto rec? \} \\
sessions' = sessions \\
claims' = claims \\
delegationOf' = delegationOf \\
epoch' = epoch
\end{schema}

Four guards earn their place against the bead's own criterion:
$\field{sid?} \notin \dom \field{claims}$ -- a session that already holds its own
claim does not need (and must not receive) a fork-propagated one, so
this operation can never silently overwrite an explicit
\field{ClaimMission}; $\field{witness?} \in \dom \field{sessions}$ together with
$\field{(sessions(witness?)).root} = \field{(sessions(sid?)).root}$ -- the
witnessing session must be a genuine, currently-rostered member of the
FORKED session's OWN lineage (the two conjuncts together are the
inline expansion of ``\field{witness?} is in \field{sid?}'s
\field{family}'', unpacked rather than named via a global function
for the reason \S\ref{sec:session-schema} gives), not an arbitrary
other session -- exactly the \field{root}/\field{family} lookup
(\S\ref{sec:harness-instantiation}) replacing a parent-chain walk; and
$\field{witness?} \in \dom \field{claims}$ -- there must be an ACTUAL claim to
propagate, not a fabricated one. No conjunct here ever needs to resolve
\field{witness?} back through \field{sessions} a second way: membership
in \field{family} already IS the traceability witness
\field{oForkPropagated} exists to name, which is exactly
\field{TraceabilityViaRootSuffices}'s claim (\S\ref{sec:theorems}).

\subsection{BindDelegationToSession}

Models \field{WriteDelegationBinding}
(\fpath{internal/mission/active.go}:349--358) -- the per-session
sidecar bridging the PreToolUse allocation point to the PostToolUse
audit writer, DES-054 v5.

\begin{schema}{BindDelegationToSession}
\Delta Registry \\
sid? : SESSIONID \\
id? : DELEGATIONID
\where
sid? \in \dom sessions \\
id? \in \dom delegations \\
(delegations(id?)).session = sid? \\
delegationOf' = delegationOf \oplus \{ sid? \mapsto id? \} \\
sessions' = sessions \\
delegations' = delegations \\
claims' = claims \\
epoch' = epoch
\end{schema}

The guard \field{(delegations(id?)).session = sid?} is exactly what
keeps \field{Registry}'s own \field{delegationOf} invariant
(\S\ref{sec:registry}) an invariant rather than a hope: a session can
only ever bind to a delegation IT ITSELF opened, never to one opened by
a different session.

\subsection{UnbindDelegation}

Models \field{ClearDelegationBinding}
(\fpath{internal/mission/active.go}:968--980).

\begin{schema}{UnbindDelegation}
\Delta Registry \\
sid? : SESSIONID
\where
sid? \in \dom delegationOf \\
delegationOf' = \{ sid? \} \ndres delegationOf \\
sessions' = sessions \\
delegations' = delegations \\
claims' = claims \\
epoch' = epoch
\end{schema}

\subsection{CorroborateSession}
\label{sec:opcorroborate}

Models re-validating a session's liveness anchor against CURRENT
reality -- \field{isLiveAncestor}'s live walk for Claude Code
(\fpath{tree.go}:136--161), or (per
\S\ref{sec:harness-instantiation}) an OpenCode session-liveness query
for OpenCode. The real OS/harness-level check this models -- is the
anchor genuinely still live -- happens OUTSIDE \field{Registry}'s own
state, the same way \field{isLiveAncestor} reads \fpath{/proc} rather
than any ethos-persisted fact; this operation's precondition captures
only the ethos-side bookkeeping requirement (the session is on record
and was not already marked dead), not the external liveness check
itself, which the caller is assumed to have already performed before
invoking it.

\begin{schema}{CorroborateSession}
\Delta Registry \\
sid? : SESSIONID \\
rec? : SESSION
\where
sid? \in \dom sessions \\
(sessions(sid?)).alive = ztrue \\
rec?.harness = (sessions(sid?)).harness \\
rec?.parent = (sessions(sid?)).parent \\
rec?.root = (sessions(sid?)).root \\
rec?.alive = ztrue \\
rec?.bornEpoch = epoch \\
sessions' = sessions \oplus \{ sid? \mapsto rec? \} \\
delegations' = delegations \\
claims' = claims \\
delegationOf' = delegationOf \\
epoch' = epoch
\end{schema}

Every field of \field{rec?} except \field{bornEpoch} is forced equal to
the session's existing record: corroboration can only refresh WHEN a
session was last confirmed live, never silently change WHAT it is.

\subsection{RestartService}
\label{sec:oprestart}

Models the class of event \field{FindClaudePID}'s own doc comment names
as the reason liveness must never be cached: a daemon restart, or any
long-lived \field{ethos serve} process moving on from the session it
last observed (\fpath{tree.go}:18--22). The real event is unconditional
-- a restart can happen at any time, and happening does not itself
discard any durable data; the \field{epoch < maxEpoch} guard below is
purely the \S\ref{sec:constants} animatability bound, with no real-code
counterpart (the real \field{epoch} has no ceiling).

\begin{schema}{RestartService}
\Delta Registry
\where
epoch < maxEpoch \\
epoch' = epoch + 1 \\
sessions' = sessions \\
delegations' = delegations \\
claims' = claims \\
delegationOf' = delegationOf
\end{schema}

Nothing here deletes a session, a claim, or a delegation binding --
every sidecar genuinely is a FILE that survives a restart on disk. What
changes is \field{epoch} alone, and that one change is sufficient:
\field{RestartClearsLiveness} (\S\ref{sec:theorems}) shows that every
session recorded before the restart immediately fails the
\field{bornEpoch = epoch} conjunct every operation in \S\ref{sec:ops}
checks, so none of them is usable again until explicitly
\field{CorroborateSession}'d.

\subsection{EndSession}
\label{sec:opendsession}

Models \field{Store.Delete}/\field{Store.Purge} removing a roster entry
and its sidecars (\fpath{internal/session/store.go}:200--301,
\field{clearMissionSidecarsLocked}) -- but NOT the delegation tree,
which is durable audit data outliving the roster (\S\ref{sec:registry}).

\begin{schema}{EndSession}
\Delta Registry \\
sid? : SESSIONID
\where
sid? \in \dom sessions \\
sessions' = \{ sid? \} \ndres sessions \\
claims' = \{ sid? \} \ndres claims \\
delegationOf' = \{ sid? \} \ndres delegationOf \\
delegations' = delegations \\
epoch' = epoch
\end{schema}

\subsection{AdapterDataMissing}
\label{sec:opadaptermissing}

The ``missing harness adapter data'' failure mode: a harness lifecycle
event fires with insufficient information to identify a session at all
-- no usable, stable ID. Rather than coerce a fabricated or
best-guess \field{SESSIONID} into existence, the event is modelled as a
self-loop that touches nothing, the same discipline
spec-mission-lifecycle.tex's \field{Idle} and
spec-writeset-admission.tex's non-goals both apply to an input the
model refuses to manufacture a value for.

\begin{schema}{AdapterDataMissing}
\Xi Registry \\
harness? : HARNESS
\end{schema}

\subsection{ResolveCommitAttribution}
\label{sec:opresolve}

Models the operation \field{ethos hook commit-trailers} performs --
\field{runHookCommitTrailers} resolving a session via
\field{resolve.SessionID}, then \field{WriteCommitTrailers} reading that
ONE session's sidecars (\fpath{cmd/ethos/hook.go}:218--251,
\fpath{internal/hook/commit\_trailers.go}:45--72) -- generalized to any
harness. A pure query: \field{Registry} is never changed by resolving a
commit's trailers.

\begin{schema}{ResolveCommitAttribution}
\Xi Registry \\
sessionExpected? : ZBOOL \\
sessionResolved? : ZBOOL \\
candidate? : SESSIONID \\
trailerMission! : MISSIONLINK \\
trailerDelegation! : DELEGLINK \\
diagnostic! : ZBOOL
\where
sessionResolved? = ztrue \implies sessionExpected? = ztrue \\
sessionExpected? = zfalse \implies
  ( trailerMission! = noMission \land trailerDelegation! = noDeleg \land
    diagnostic! = zfalse ) \\
( sessionExpected? = ztrue \land sessionResolved? = zfalse ) \implies
  ( trailerMission! = noMission \land trailerDelegation! = noDeleg \land
    diagnostic! = ztrue ) \\
( sessionResolved? = ztrue \land
  \lnot ( candidate? \in \dom sessions \land
          (sessions(candidate?)).alive = ztrue \land
          (sessions(candidate?)).bornEpoch = epoch \land
          candidate? \in \dom claims ) ) \implies
  ( trailerMission! = noMission \land trailerDelegation! = noDeleg \land
    diagnostic! = zfalse ) \\
( sessionResolved? = ztrue \land
  candidate? \in \dom sessions \land
  (sessions(candidate?)).alive = ztrue \land
  (sessions(candidate?)).bornEpoch = epoch \land
  candidate? \in \dom claims ) \implies
  ( trailerMission! = missionOf(claims(candidate?)) \land
    diagnostic! = zfalse \land
    ( ( candidate? \in \dom delegationOf \land
        (delegations(delegationOf(candidate?))).mission = missionOf(claims(candidate?)) )
      \implies trailerDelegation! = delegOf(delegationOf(candidate?)) ) \land
    ( \lnot ( candidate? \in \dom delegationOf \land
              (delegations(delegationOf(candidate?))).mission = missionOf(claims(candidate?)) ) )
      \implies trailerDelegation! = noDeleg )
\end{schema}

Four cases, mutually exclusive on \field{sessionExpected?}/
\field{sessionResolved?}/the candidate's own usability, and each one
answers a question the bead's acceptance criteria ask directly:

\begin{itemize}
  \item \textbf{No session was ever expected} (headless, CI, a plain
    terminal, or -- generalized -- any invocation outside either
    harness): silent, no diagnostic. This is \field{UnderClaudeCode()}
    returning false and \field{SessionID} returning
    \field{ErrNotUnderClaudeCode}
    (\fpath{resolve.go}:233--247) -- a normal state, not a failure.
  \item \textbf{A session was expected but could not be resolved} --
    ambiguous attribution, a broken pointer, an uncorroborated anchor:
    no trailer, AND \field{diagnostic!} = \field{ztrue}. This is
    \field{SessionID} returning \field{ErrNoSession}
    (\fpath{resolve.go}:154--197), and \field{runHookCommitTrailers}'s
    own stderr line (\fpath{hook.go}:237--239) is exactly this
    conjunct's observable effect: ``\emph{a silently missing trailer is
    exactly the failure class it exists to prevent}''
    (\fpath{hook.go}:230--233).
  \item \textbf{A session resolved cleanly but has no usable claim} --
    dead, uncorroborated since the last restart, or simply never
    claimed anything: no trailer, and -- the point that must not be
    conflated with the previous case -- NO diagnostic either. This is
    the bead's own ``a small, mission-free leader commit never receives
    a Mission/Delegation trailer'' criterion, and the silence is
    deliberate: an unclaimed session committing is the ORDINARY case,
    not a resolution failure, exactly matching
    \field{WriteCommitTrailers}'s own silent return when
    \field{missionID == ""} (\fpath{commit\_trailers.go}:53--55).
  \item \textbf{A session resolved and has a usable claim}: the mission
    trailer is emitted, and the delegation trailer is emitted ONLY when
    the session's own delegation-binding sidecar names a delegation
    bound to the SAME mission -- precisely
    \field{WriteCommitTrailers}'s own
    \field{b.MissionID != missionID} gate
    (\fpath{commit\_trailers.go}:65), restated here as a direct
    consequence of \field{Registry}'s own \field{delegationOf}
    invariant (\S\ref{sec:registry}) rather than a read-time check this
    operation has to perform itself.
\end{itemize}

\field{FailClosedOnAmbiguity} (\S\ref{sec:theorems}) names the first
two cases' distinction as the theorem it is.

% ===================================================================
\section{Theorems}
\label{sec:theorems}
% ===================================================================

\fuzz{} type-checks schemas; it does not discharge proof obligations.
The following are claims about the operation set
$\{\field{JoinRootSession}, \field{ForkSession}, \field{ClaimMission},
\field{ReleaseClaim}, \field{AllocateDelegationTierA},
\field{AllocateDelegationTierB}, \field{AllocateDelegationForkPropagated},
\field{BindDelegationToSession}, \field{UnbindDelegation},
\field{CorroborateSession}, \field{RestartService}, \field{EndSession}\}$,
each backed by inspection of the schemas above; two
(\field{DelegationIDFreshness}, \field{RestartClearsLiveness}) are
additionally exercised by ProB animation (\S\ref{sec:validation}). None
are discharged by proof.

\begin{center}
\begin{tabularx}{\textwidth}{@{}lX@{}}
\toprule
\textbf{Theorem} & \textbf{Justification} \\
\midrule
\field{DelegationIDFreshness} &
  \field{AllocateDelegationTierA}, \field{AllocateDelegationTierB}, and
  \field{AllocateDelegationForkPropagated} all guard
  $\field{id?} \notin \dom \field{delegations}$, and no OTHER operation in the set
  adds to $\dom \field{delegations}$ at all. Two concurrent
  sessions -- same or different harness -- can therefore never end up
  bound to the same \field{delegation\_id}: whichever atomic step
  allocates first makes the id unavailable to every subsequent one,
  directly answering the bead's first acceptance criterion. Real
  mutual exclusion across truly concurrent allocators is provided by
  \field{AcquireDispatchPendingLock} (\S\ref{sec:nongoals}'s disclosed
  scope), not re-derived here. \\
\field{ForkNeverInheritsClaim} &
  \field{JoinRootSession}'s and \field{ForkSession}'s own frame
  conjunct \field{claims' = claims} (\S\ref{sec:opjoin},
  \S\ref{sec:opforksession}) holds identically for a root session and a
  forked one -- there is no case split in either schema's guard that
  could let a fork take a different path to a claim than a root session
  does. A freshly joined OR forked session therefore has NO
  \field{claims} entry the instant it enters the registry; any later
  entry requires a separate, explicit \field{ClaimMission}, an
  \field{AllocateDelegationTierB} whose \field{oInherited} guard already
  forces a recorded, checkable \field{parentDelegation} link
  (\S\ref{sec:delegation-schema}), or an
  \field{AllocateDelegationForkPropagated} whose guard requires an
  ACTUAL witnessing claim on a genuine \field{family} member
  (\S\ref{sec:opforkpropagate}). Silent inheritance is not merely
  guarded against; it has no value in \field{ORIGIN} that could express
  it without one of those two checkable witnesses. Confirming that V2
  fork is frequent and first-class (\S\ref{sec:harness-instantiation})
  sharpens this claim rather than weakening it: \field{ForkSession} is
  not a rarely-exercised edge case whose guard might go untested in
  practice, but the routine path every multi-step OpenCode mission round
  is expected to take. \\
\field{TraceabilityViaRootSuffices} &
  \field{AllocateDelegationForkPropagated}'s guard
  (\S\ref{sec:opforkpropagate}) requires \field{witness?} to be a member
  of \field{sid?}'s family -- which the schema states, unpacked, as a
  single \field{root} equality check against \field{witness?}, needing
  no traversal of any chain. Because
  \field{ForkSession} stamps a correct, inherited \field{root} on every
  session it creates (its own guard,
  \field{rec?.root = (sessions(parentSid?)).root}), and
  \field{JoinRootSession} stamps every root session as its own root,
  \field{root} equality between two live sessions is ALWAYS sufficient
  evidence that they share a lineage -- there is no additional
  provenance fact \field{AllocateDelegationForkPropagated} needs to
  consult beyond what \field{family} already packages. This is the
  formal statement of why this document's second draft could replace a
  parent-chain WALK (the first draft's rejected design) with a single
  \field{root} LOOKUP once OpenCode's own API was confirmed to expose
  lineage the same way. \\
\field{RestartClearsLiveness} &
  \field{RestartService}'s sole effect is \field{epoch' = epoch + 1}; no
  session's \field{bornEpoch} changes. Every operation in
  \S\ref{sec:ops} that reads a session's liveness checks
  \field{(sessions(s)).bornEpoch = epoch} as one of its two liveness
  conjuncts, so the instant a restart fires, every PRE-restart session
  fails that check and is unusable by any operation until
  \field{CorroborateSession} explicitly re-stamps it. A stale on-disk
  binding is therefore structurally incapable of being treated as live
  across a restart without a fresh, explicit corroboration step --
  exactly \field{FindClaudePID}'s own ``\emph{resolved fresh on every
  call -- never cached}'' discipline (\fpath{tree.go}:18--22), lifted
  from one session's PID to every session's epoch. \\
\field{WorktreeIndependence} &
  Every $\Delta \field{Registry}$ operation above updates
  \field{claims}/\field{delegationOf}/\field{sessions} via
  $\oplus$ or $\ndres$ applied to EXACTLY ONE \field{sid?}
  key, never a blanket update over a range or a second session's
  entry. Two distinct sessions -- whether or not either is bound to a
  worktree (this document has no separate \field{WORKTREE} type,
  because nothing in the model needs to single it out: the
  independence property holds for ANY two distinct sessions, a strictly
  stronger and more economical result than modelling worktree as a
  special case) -- can therefore never have one operation's effect on
  one cross into the other's state. \\
\field{HarnessNeutrality} &
  No guard or effect predicate in any schema in \S\ref{sec:ops}
  mentions the \field{harness} field at all. The entire state machine
  -- and hence every invariant and theorem in this table -- therefore
  applies identically to a session regardless of which \field{HARNESS}
  value it carries: \S\ref{sec:modelling-choice}'s design choice (one
  registry, not two) is exactly what makes this claim checkable by
  inspection rather than asserted by analogy. \\
\field{FailClosedOnAmbiguity} &
  \field{ResolveCommitAttribution}'s second and third cases
  (\S\ref{sec:opresolve}) both emit no trailer, but differ on
  \field{diagnostic!}: \field{ztrue} exactly when a session was
  expected and did not resolve (an ambiguous or broken binding -- the
  failure class the bead's criterion names), \field{zfalse} when a
  session resolved cleanly but simply has no claim (the ordinary,
  unremarkable mission-free-commit case). The schema's own guard
  structure makes these two outcomes mutually exclusive and jointly
  exhaustive of ``no trailer was emitted,'' so a mission-free leader
  commit and a genuinely ambiguous one can never be confused for each
  other at the output. \\
\bottomrule
\end{tabularx}
\end{center}

% ===================================================================
\section{Harness Instantiation}
\label{sec:harness-instantiation}
% ===================================================================

\textbf{On this section's sourcing.} This agent's own toolset for the
present mission carried no live web-fetch or web-search capability.
The OpenCode-side mapping below is therefore attributed, throughout, to
a companion research artifact
(\fpath{docs/prfaq-opencode-integration.md}, a sibling worktree) that
WAS produced with direct source access -- fetched via context7 from
\fpath{anomalyco/opencode}'s \field{dev} branch across two evaluation
rounds, the second of which both corrected a V1-vocabulary draft and
located the \field{context.data.session}
hierarchy/liveness API a first search had correctly disclosed as
unfound rather than guessed at. OpenCode V2 is a RELEASED, currently
running system (v2.0.26 confirmed installed, in production use on
multiple operator machines) -- an earlier round of this same research
mistakenly reported V2 as dev-branch-only and not yet cut; that was
wrong, was retracted on the research mission's own audit trail, and is
corrected here: every claim below targets the CURRENT, shipping V2
system, not a forthcoming one. Every quoted or paraphrased OpenCode
fact below carries its own citation back to the specific section of the
companion artifact (or, transitively, the companion artifact's own
primary source) it is attributed to; this author did not independently
re-fetch any of those primary sources. The Claude Code mapping that
follows carries no such attribution chain: it is read directly off this
repository's own code, cited by file and line, exactly as every other
Claude-Code-side claim in this document.

\subsection{Claude Code}

\fpath{SESSION.harness} = \field{hClaude}. \fpath{SESSIONID} is a
corroborated \field{CLAUDE\_PID}
(\field{ClaudePIDFromEnvCorroborated}, \fpath{tree.go}:68--86) or an
explicit \field{ETHOS\_SESSION} value (\fpath{resolve.go}:316--318),
whichever \field{SessionID}'s chain resolves first.

The striking fact, confirmed directly from \fpath{tree.go}'s own doc
comment (lines 28--33): \textbf{\field{SESSION.parent} is, in every real
Claude Code binding today, \field{noParent}}. ``\emph{a Task-tool
subagent or hook subprocess inherits its leader's \field{CLAUDE\_PID}
verbatim, and correctly so, since it belongs to the same session}'' --
Claude Code never mints a new top-level \field{SESSIONID} for a
subagent at all; leader and subagent share one session, full stop. The
forking Claude Code DOES have lives entirely at the DELEGATION level --
\field{parent\_delegation}, already modelled by
\field{AllocateDelegationTierB}'s \field{oInherited} case
(\S\ref{sec:opallocateb}) -- not at the session level this document's
\field{PARENTLINK} field exists to generalize. \field{ForkNeverInheritsClaim}
is, for Claude Code specifically, true in a nearly vacuous sense: the
session-level fork it protects against essentially never occurs in the
first place, because Claude Code's own architecture never creates the
precondition.

\field{IsLiveSession}'s two conjuncts instantiate directly:
\field{alive} = whether the session's roster file is still present
(\field{Store.Load} succeeds), and \field{bornEpoch = epoch} =
\field{isLiveAncestor}'s own fresh-every-call re-walk of
\fpath{/proc} (or the macOS/Windows equivalent) confirming the
corroborated PID is STILL a live ancestor right now, not merely was one
when first observed.

\subsection{OpenCode}
\label{sec:opencode-instantiation}

\fpath{SESSION.harness} = \field{hOpenCode}. An ethos-owned plugin is a
\field{Plugin.define(\{id, async setup(ctx) \{...\}\})} module, loaded
from \fpath{.opencode/plugins/} (or an npm package named in
\fpath{opencode.json}'s \field{"plugins"} array), whose \field{setup}
subscribes via \field{ctx.event.subscribe(\{signal\})} -- consumed with
\field{for await} inside \field{setup} -- to session lifecycle event
types including \field{session.created}, \field{session.updated},
\field{session.idle}, and \field{session.deleted}
(\fpath{docs/prfaq-opencode-integration.md}, \S~``Why not just reuse
the Claude Code hook mechanism verbatim?,'' citing
\fpath{anomalyco/opencode}'s
\fpath{services/www/src/docs/content/build/plugins/migrate-v1.mdx},
confirmed via context7). \field{ctx.session.hook()} is a DIFFERENT,
narrower surface accepting only \field{"context"}/\field{"prompt"}/
\field{"compaction"}/\field{"model.request"}/\field{"http.request"} --
none of which fire on session start or update -- so
\field{ctx.event.subscribe()} is the confirmed, correct primitive for
this integration's \field{JoinRootSession}/\field{ForkSession}
triggers, identified by elimination against that narrower list in the
companion research.

\textbf{JoinRootSession and ForkSession.} On a \field{session.created}
event with no fork ancestry, the plugin calls \field{ethos mission
claim <id>} (the existing Tier B dispatch-binding primitive, per this
repo's own \fpath{docs/development.md} Build \& Run section) if an
active claim is relevant, instantiating \field{JoinRootSession}
(\S\ref{sec:opjoin}). A session FORK is OpenCode's own
\field{POST /api/session/\{sessionID\}/fork} HTTP endpoint (operation ID
\field{session.fork}, confirmed listed in
\fpath{anomalyco/opencode}'s \fpath{V2\_HTTP\_API\_AUDIT.md} on the
\field{dev} branch, Group 4 table, row 053) -- the forked child gets its
OWN, distinct OpenCode session ID; the V2 session contract spec states
directly that forking ``copies messages only through its selected
boundary but adopts the parent's newest instruction values as its
baseline'' (\fpath{specs/v2/session.md}, \S~``Instructions Are Value
Deltas''), and nothing in that contract ties the two session identities
together beyond transcript and instruction-value copying. This
instantiates \field{ForkSession} (\S\ref{sec:opforksession}) directly:
\field{session.fork}'s own confirmed semantics are exactly why this
document's \field{rec?.root} guard inherits from the PARENT rather than
starting a new lineage, and exactly why \field{claims' = claims} is the
one frame conjunct that must hold regardless.

\textbf{Lineage lookup: root/family as confirmed, maintained
projections, not payload fields.} The CLI-plugin data-cache API --
\field{context.data.session.\{list, get, root, family, cost, status,
sync, invalidate\}}, on a distinct plugin surface
(\fpath{@opencode/plugin/tui}, loaded via a package's \field{./tui}
export) that coexists alongside the main hook-plugin entry in the same
package -- exposes \field{context.data.session.root(sessionID)} and
\field{.family(sessionID)} as DIRECT hierarchy lookups, confirmed by a
full read of \fpath{anomalyco/opencode}'s
\fpath{services/www/src/docs/content/build/plugins/cli.mdx} on the
\field{dev} branch (\fpath{docs/prfaq-opencode-integration.md},
\S~``What happens if OpenCode forks a session mid-mission?''). This is
the exact confirmation that licenses \S\ref{sec:session-schema}'s own
design choice: \field{SESSION.root} is stored, maintained data, not
something a plugin -- or this document's own
\field{AllocateDelegationForkPropagated}, \S\ref{sec:opforkpropagate}
-- has to reconstruct by chaining \field{session.created} event payload
fields together, whose exact field set the companion research explicitly
leaves at medium confidence (\fpath{prfaq-opencode-integration.md},
\S~``What I could not verify''). \field{TraceabilityViaRootSuffices}
(\S\ref{sec:theorems}) is the formal payoff of this confirmation.

\textbf{Liveness: \field{context.data.session.status}, not a
reconstructed execution-claim check.} OpenCode's V2 session contract
commits a write-ahead claim per process-local busy period, preserved
across an unclean server restart but explicitly NOT a liveness
guarantee -- ``\emph{durable execution events and claims are historical
and recovery records, not proof that this process is still live}''
(\fpath{specs/v2/session.md}, \S~``Execution Is Process-Local,'' as
quoted in the companion research). Reconstructing \field{IsLiveSession}
from that recovery language directly would be exactly the expensive,
indirect check the companion research's restart FAQ warns against;
\field{context.data.session.status(sessionID)} is confirmed as a
DIRECT, maintained liveness projection instead -- materially cheaper
than polling \field{sessions.active()} and cross-referencing
recovery-attempt counts. This document's \field{SESSION.alive}
instantiates to \field{status(sessionID)} for OpenCode, the same
abstract role \field{Store.Load} succeeding plays for Claude Code
(below); \field{SESSION.bornEpoch = epoch} instantiates to a fresh
\field{status} query returning ``live'' RIGHT NOW, not merely having
returned it once.

\textbf{Worktree keying.} V2 replaces V1's flat \field{directory}
context field with \field{ctx.location.directory}; a separate, flat
\field{worktree} field was removed entirely in favor of a structured
\field{ctx.worktree.\{create,list,refresh,remove\}} domain API keyed by
\field{projectID}. For this integration's correlation-keying purpose,
\field{ctx.location.directory} alone is sufficient: a git worktree
checkout IS a directory, and \field{.punt-labs/ethos/} is a tracked
directory present identically inside every worktree, so the plugin
keys its write by \field{(ctx.location.directory, OpenCode Session
ID)} with no new multi-worktree design needed
(\fpath{docs/prfaq-opencode-integration.md}, \S~``What happens with
OpenCode worktree or concurrent sessions?''). This grounds
\field{WorktreeIndependence} (\S\ref{sec:theorems}) for OpenCode the
same way the existing per-worktree \fpath{.punt-labs/ethos/} layout
already grounds it for Claude Code.

The \field{AdapterDataMissing} failure mode
(\S\ref{sec:opadaptermissing}) is named for exactly this harness: an
OpenCode plugin lifecycle event firing with a payload that does not
carry a usable, stable session identifier (a malformed or partial
event, a hook firing before the server has finished assigning an ID) is
the shape of failure this operation's self-loop represents, where
Claude Code's own \field{CLAUDE\_PID}-corroboration chain has no direct
analogue to fail in quite the same way.

% ===================================================================
\section{Failure Modes}
\label{sec:failure-modes}
% ===================================================================

The task brief names three; each is addressed by a construct already
introduced above, gathered here for a single point of reference.

\begin{itemize}
  \item \textbf{Ambiguous attribution.}
    \field{ResolveCommitAttribution}'s second case
    (\S\ref{sec:opresolve}): $\field{sessionExpected?} = \field{ztrue} \land
    \field{sessionResolved?} = \field{zfalse}$ emits no trailer and sets
    $\field{diagnostic!} = \field{ztrue}$. \field{FailClosedOnAmbiguity}
    (\S\ref{sec:theorems}) is the formal statement that this case can
    never be confused with the silent, undiagnosed mission-free-commit
    case.
  \item \textbf{Missing harness adapter data.} \field{AdapterDataMissing}
    (\S\ref{sec:opadaptermissing}): a lifecycle event with no usable
    session identifier is modelled as a self-loop touching nothing in
    \field{Registry}, rather than coercing a fabricated
    \field{SESSIONID} into existence. See
    \S\ref{sec:harness-instantiation} for why this is named for
    OpenCode specifically.
  \item \textbf{Concurrent writers to the same sidecar.}
  \label{sec:failure-concurrent}
    Disclosed as out of scope in \S\ref{sec:nongoals}: this is a Z
    document of atomic, sequential-step operations, the identical
    modelling boundary spec-mission-lifecycle.tex and
    spec-writeset-admission.tex both disclose for their own state
    machines. What IS checked here is the structural property that
    makes the real locks (\field{AcquireDispatchPendingLock}, the
    per-session roster flock) SUFFICIENT to prevent corruption once
    they serialize these operations: every $\Delta \field{Registry}$
    operation above writes to \field{claims}/\field{delegationOf}/
    \field{sessions}/\field{delegations} via a SINGLE-KEY $\oplus$
    or $\ndres$, never a multi-key sweep -- so SERIALIZED
    execution of these operations (which the real flocks guarantee, and
    this document does not re-derive) composes safely, by the same
    reasoning \field{WorktreeIndependence} already gives.
\end{itemize}

% ===================================================================
\section{Validation}
\label{sec:validation}
% ===================================================================

\texttt{/z-spec:check} (\texttt{fuzz}) type-checks this document with
\textbf{0 errors}.

\subsection{Certified-complete run at setsize~1}

\texttt{probcli -model\_check -p DEFAULT\_SETSIZE 1 -df -c} (raising
\field{MAX\_OPERATIONS} from its default of 10, which otherwise
truncates transition computation before the state space is actually
exhausted) explores \field{SESSIONID}, \field{MISSIONID}, and
\field{DELEGATIONID} each at a single witness value and reports
\textbf{``No counter example found. ALL states visited'' -- 287
states, 3862 transitions, certified complete}: \textbf{zero invariant
violations, zero deadlocks}, over the ENTIRE reachable state space at
this setsize, not a bounded sample of it. Sixteen of this document's
eighteen operations fire -- \field{ResolveCommitAttribution} alone
1347 times, \field{AdapterDataMissing} 570 times, every claim/release/
corroborate/bind/unbind/end transition at least a dozen times each.
\field{ForkSession} and \field{AllocateDelegationForkPropagated} are
the two operations the coverage census reports as
\texttt{UNCOVERED\_OPERATIONS} at this setsize, and the reason is
structural, not a defect: both require TWO DISTINCT \field{SESSIONID}
values in their own guards (\field{ForkSession}'s
$\field{sid?} \notin \dom \field{sessions} \land \field{parentSid?} \in \dom \field{sessions}$,
\field{AllocateDelegationForkPropagated}'s
$\field{witness?} \in \dom \field{sessions}$ naming a DIFFERENT session from
\field{sid?}), which a carrier of size~1 cannot supply -- the
coverage gap is forced by the setsize bound alone, confirmed by its
disappearance the moment setsize rises to 2 (below), not a sign either
operation is unreachable in general.

\subsection{Bounded run at setsize~2}

At \field{DEFAULT\_SETSIZE} 2 ($\field{maxEpoch} = 2$,
$\field{maxSessions} = 2$, $\field{maxDelegations} = 2$, per
\S\ref{sec:constants}) -- the setsize every other check in this
document's own toolchain defaults to, and the one at which
\field{ForkSession} finally has the second \field{SESSIONID} value its
own guard needs -- the reachable space is considerably LARGER than
either existing specification's own (spec-mission-lifecycle.tex: 373
states; spec-writeset-admission.tex: 5790 states).

\textbf{Errata (second correction): probcli's own transition count at
this setsize is NOT uniformly non-deterministic, and the first
correction to this section misattributed its cause.} That correction
observed three different transition counts under nominally the same
\texttt{-mc 2000 -df} invocation -- 49,062, 82,155, 50,400 -- across
three independent invocations by three parties, and attributed the
split to ``probcli's own search order.'' This is WRONG, and is itself
corrected here. The evaluator's final-round reproduction, independently
confirmed by rerunning both exact invocations here, isolates the real
cause: \texttt{-p MAX\_OPERATIONS 5000 -mc 2000 -df} returns
\textbf{82,155} transitions, deterministically, every time (3 runs
confirmed by the evaluator, 3 more confirmed here); \texttt{-mc 2000
-df} with probcli's DEFAULT \field{MAX\_OPERATIONS} returns
\textbf{50,400}, deterministically, every time (2 runs confirmed by the
evaluator, 2 more confirmed here). The 82,155/50,400 split is a
\textbf{\field{MAX\_OPERATIONS} CONFOUND between two differently-flagged
invocations across the first and second drafts of this section} -- NOT
search-order non-determinism, and NOT evidence that \texttt{-df} itself
is non-deterministic. Genuine run-to-run non-determinism IS real, but
is confirmed ONLY under the no-flags/\texttt{-c}-alone configuration:
the evaluator observed five DIFFERENT counts across five runs of
\texttt{-mc 2000} with no \texttt{-df} (47,254, 48,111, 47,158, 46,780,
47,852), and this document's own earlier reruns independently observed
the same pattern (ten different counts across ten no-flags/\texttt{-c}
runs, a two-run spread at a 10,000-state no-flags bound). The original
49,062 figure this whole errata chain traces back to was never
reproduced by anyone under either deterministic \texttt{-df}
configuration, and its exact provenance (which \field{MAX\_OPERATIONS}
value, if any, the very first invocation used) is not reconstructible
from the record; it does not need to be, since the corrected, properly
attributed picture below supersedes it.

\textbf{Errata (second correction): the claim that all eighteen
operations fire ``under every flag configuration tried'' is false, and
is corrected here.} \field{AllocateDelegationForkPropagated} is
UNCOVERED in every \texttt{-df} run performed by the evaluator (5 runs,
both \field{MAX\_OPERATIONS} variants) and independently reconfirmed
here (3 further runs, both variants): \textbf{18/18 coverage is
achieved ONLY under the no-flags configuration}, not under
\texttt{-df}. This was checked directly, not inferred: the
\texttt{probcli -c} coverage census under \texttt{-mc 2000 -df -c}
(both with and without \field{MAX\_OPERATIONS} raised to 5000) reports
\texttt{UNCOVERED\_OPERATIONS (1), AllocateDelegationForkPropagated} in
every run; the identical invocation with \texttt{-df} dropped reports
\texttt{UNCOVERED\_OPERATIONS (0)} every time this document's own
reruns tried it. The earlier draft's claim of universal coverage
``under every flag configuration tried'' did not hold even at the time
it was written, because \texttt{-df}'s own deadlock-search discipline
evidently visits states in an order that starves this one operation's
guard of a satisfying witness within the explored bound -- a fact about
\texttt{-df}'s search strategy, not about this specification's
reachability (\S\ref{sec:opforkpropagate}'s own guard is satisfiable,
confirmed directly by its firing under the no-flags configuration at
the identical 2000-state bound).

What DOES reproduce, consistently, across EVERY configuration and every
run performed by all three parties, with no exception found: \textbf{zero
invariant violations, zero deadlocks}, at bounds ranging from 2000
states up to 10,000+ states (the evaluator's own independent check:
414,730 transitions, 0 violations) and 13,632 states (this document's
own earlier, longer run: 317,844 transitions, 0 violations). This
safety claim is NOT weakened by either correction above -- it is the
one claim in this subsection that survives both.

This IS a disclosed, not a hidden, coverage gap at setsize~2 -- the
exact discipline spec-mission-lifecycle.tex's own
\S~``A disclosed, pre-existing gap'' models, now corrected TWICE over:
first, in an earlier revision, for treating a flag-dependent number as
a fixed fact and building a withdrawn growth-rate narrative on it; and
second, here, for misattributing the flag-dependence itself to search
non-determinism and for overstating operation coverage. What IS
established without
qualification: the ENTIRE setsize~1 state space (287 states) is free of
invariant violations and deadlocks, with sixteen of eighteen operations
exercised and the remaining two's non-firing explained structurally;
zero invariant violations and zero deadlocks hold across every
setsize~2 run performed by any party under any configuration, with no
exception found; and full 18/18 operation coverage is achieved under
the no-flags configuration specifically, confirmed at every bound
tried from 2000 through 10,000+ states. What is NOT established: the
exact size of the setsize~2 reachable space; a single, configuration-
independent transition count for any given state bound (two of the
three counts this errata chain discusses are flag-determined, not
run-to-run noise -- only the no-flags/\texttt{-c}-alone configuration
is genuinely non-deterministic); any claim about how a transition count
trends as the state bound grows (withdrawn in the first correction and
not restored here); or full operation coverage under \texttt{-df}
specifically, which this correction now states does not hold.
Diagnosing why \texttt{-df}'s own search order leaves
\field{AllocateDelegationForkPropagated} unreached within a 2000-state
bound, and whether a larger bound or different flags close that gap,
is scoped as follow-up work for a subsequent revision, not resolved
here.

% ===================================================================
\section{Open Dependency: punt-labs/points}
\label{sec:open-dependency}
% ===================================================================

Bead \field{ethos-5y4r} names a mirrored bead, \field{points-hqi}, in
\fpath{punt-labs/points} -- the repo responsible for configuring how AI
coding harnesses route to model sources, and therefore a plausible
owner of harness-adapter-side work this bead's acceptance criteria
depend on. A search of that repository's issue tracker via \field{gh}
for ``\field{points-hqi}'', ``\field{ethos}'', and ``\field{session
traceability}'' during this drafting session returned ZERO results --
either the issue is tracked elsewhere, the tracker is not populated for
this bead yet, or issues are disabled on that repository. This is
recorded here, honestly, as an UNRESOLVED open dependency: the content
of \field{points-hqi} could not be read and is not reflected anywhere
in this document. Confirming its content (or its absence) is a
prerequisite for the next pipeline stage's design review, not something
this specification can discharge on its own.

% ===================================================================
\section{Traceability}
\label{sec:trace}
% ===================================================================

\begin{center}
\begin{tabularx}{\textwidth}{@{}llX@{}}
\toprule
\textbf{Z construct} & \textbf{Go source} & \textbf{Note} \\
\midrule
\field{HARNESS} & -- & new; names the two harnesses this bead covers \\
\field{TIER} & \fpath{internal/mission/delegation.go}:19--20 & \\
\field{ORIGIN} & \field{BindOrigin*}/\field{BoundVia*}
  (\fpath{active.go}:100--113, \fpath{delegation.go}:59--62) &
  \field{oClaim}/\field{oDispatch}/\field{oInherited} collapse the
  existing Claude Code constants; \field{oForkPropagated} is new, see
  \S\ref{sec:opforkpropagate} \\
\field{SESSION} & \fpath{internal/session/session.go}'s
  \field{Participant}/\field{Roster} $+$ liveness from
  \fpath{internal/process/tree.go} & the generalized harness-neutral
  shape; \field{root} new in this redesign, modelling
  \field{context.data.session.root}, see
  \S\ref{sec:harness-instantiation} for both instantiations \\
\field{DELEGATION} & \fpath{internal/mission/delegation.go}'s
  \field{Delegation} record (fields \field{Session}, \field{Mission},
  \field{Tier}, \field{BoundVia}, \field{ParentDelegation}) & \\
\field{Registry} & no single Go type -- the session roster directory
  plus the mission tree's delegation records, read jointly & see
  \S\ref{sec:registry}'s own note on why a delegation's session field
  is NOT required to stay in $\dom \field{sessions}$ \\
``family'' (inline expansion, no Z declaration) &
  \field{context.data.session.family} & new concept; unpacked at its
  one call site as a \field{root}-equality check, not declared as a
  global function, see \S\ref{sec:session-schema} \\
\field{JoinRootSession} & \field{Store.CreateInCheckout}
  (\fpath{internal/session/store.go}:53--94); OpenCode's
  \field{session.created} event with no fork ancestry & \\
\field{ForkSession} & OpenCode's \field{POST /api/session/\{id\}/fork}
  (\field{session.fork}) & new; no Claude Code analogue --
  \S\ref{sec:harness-instantiation} shows Claude Code never exercises
  this operation \\
\field{ClaimMission} & \field{WriteActiveMission}
  (\fpath{active.go}:223--229) & \\
\field{ReleaseClaim} & \field{ClearActiveMission}
  (\fpath{active.go}:414--423) & \\
\field{AllocateDelegationTierA} & \field{dispatchTierA}
  (\fpath{internal/hook/pretooluse\_dispatch.go}:657--700) & \\
\field{AllocateDelegationTierB} & \field{dispatchTierB}
  (\fpath{pretooluse\_dispatch.go}:711--912) & \\
\field{AllocateDelegationForkPropagated} & the companion research's
  recommended OpenCode plugin behavior on a forked
  \field{session.created} & new; no Claude Code analogue, see
  \S\ref{sec:opforkpropagate} \\
\field{BindDelegationToSession} & \field{WriteDelegationBinding}
  (\fpath{active.go}:349--358) & \\
\field{UnbindDelegation} & \field{ClearDelegationBinding}
  (\fpath{active.go}:968--980) & \\
\field{CorroborateSession} & \field{isLiveAncestor}
  (\fpath{tree.go}:136--161); OpenCode's
  \field{context.data.session.status(sessionID)} & \\
\field{RestartService} & the event class \field{FindClaudePID}'s doc
  comment names (\fpath{tree.go}:18--22, DES-074) & no single
  call site -- a daemon restart or a long-lived \field{ethos serve}
  outliving an observed session \\
\field{EndSession} & \field{Store.Delete} via
  \field{deleteFiles}/\field{clearMissionSidecarsLocked}
  (\fpath{internal/session/store.go}:200--348) & \\
\field{AdapterDataMissing} & -- & new; named for a malformed OpenCode
  lifecycle payload, see \S\ref{sec:harness-instantiation} \\
\field{ResolveCommitAttribution} & \field{runHookCommitTrailers}
  (\fpath{cmd/ethos/hook.go}:218--251) via \field{resolve.SessionID}
  (\fpath{internal/resolve/resolve.go}:278--360) and
  \field{WriteCommitTrailers}
  (\fpath{internal/hook/commit\_trailers.go}:45--72) & \\
\bottomrule
\end{tabularx}
\end{center}

\end{document}
